%% kapa: exam=1 round=28 year=2017 period=1 form=A subject=civil_law subject_ko=민법 seq=34 number=34
%% source: Q-Net 감정평가사 기출문제 (cst00309) artlSeq=5199603 file=extracted/1008/5199603/2017년 제28회 감정평가사 1차시험 문제지/1교시 A형.pdf
%% conversion: kordoc 4.24.51 -> tools/gen1.py (auto); review=proofread-sonnet; flags=none
\begin{question}{34}
지역권에 관한 설명으로 옳지 않은 것은? (다툼이 있으면 판례에 따름)
\begin{choices}
  \item 1필의 토지 일부를 승역지로 하여 지역권을 설정할 수 있다.
  \item 요역지가 공유인 경우 요역지의 공유자 1인이 지역권을 취득하면 다른 공유자도 이를 취득한다.
  \item 지역권은 요역지와 분리하여 양도하지 못한다.
  \item 승역지 소유자는 지역권에 필요한 부분의 토지소유권을 지역권자에게 위기(委棄)함으로써 지역권행사를 위하여 계약상 부담하는 공작물 수선의무를 면할 수 있다.
  \item 다른 특별한 사정이 없다면 통행지역권을 시효취득한 자는 승역지 소유자가 입은 손해를 보상하지 않아도 된다.
\end{choices}
\end{question}
\answer{5}
